\documentclass[11pt]{article}
\usepackage[margin=1in]{geometry}
\usepackage{booktabs}
\usepackage{longtable}
\usepackage{hyperref}
\usepackage{enumitem}
\title{Independent Replication Report\\
\large ``Nanopore sequencing for fast determination of plasmids, phages,
virulence markers, and antimicrobial resistance genes in Shiga
toxin-producing \textit{E.\ coli}''\\
Gonz\'alez-Escalona et al., \textit{PLoS ONE} \textbf{14}(7):e0220494 (2019)}
\author{Ollie (OpenClaw AI) --- BVBRC-100 Replication Wave, target BVBRC-36}
\date{Report: 2026-07-01 \quad Backfilled to 8-artifact standard: 2026-07-05}
\begin{document}
\maketitle

\section{Paper summary}
Gonz\'alez-Escalona et al.\ closed three STEC O26:H11/H- genomes (CFSAN027343,
CFSAN027346, CFSAN027350) using MinION nanopore + PacBio long reads and compared them to
short-read Nextera-XT MiSeq. The central methodological thesis: long-read sequencing
recovers plasmids, Stx phages, virulence genes, and AMR genes that short-read assembly
fails to detect or misassembles, because those elements are mobile/repeat-rich and
fragment under short-read assembly. Data-grounded claims include closed chromosomes
$\sim$5.7/5.6/5.4 Mb, plasmid inventories (88 kb; 95+72 kb; 157 kb), MLST ST21/ST21/ST29,
strain-specific virulome patterns (Table 7), Stx phage types (Table 8), and an AMR
inventory in which only CFSAN027346 carries acquired resistance genes, all on its
$\sim$72 kb plasmid. Deposited: GenBank CP037941--CP037947.

\section{Claims tested}
\begin{longtable}{clcc}
\toprule
\# & Claim & Type & Tested \\
\midrule
C1 & Closed genomes + plasmid replicons publicly deposited & Data & Yes (7/7 pulled) \\
C2 & Chromosome sizes 5.7/5.6/5.4 Mb & Genome & Yes \\
C3 & Plasmid inventory: 343 $\to$1(88kb); 346$\to$2(95+72); 350$\to$1(157) & Genome & Yes \\
C4 & MLST 343=ST21, 346=ST21, 350=ST29 & Typing & Yes \\
C5 & AMR only in 346 (six genes on 72kb plasmid) & AMR & Yes (EXACT) \\
C6 & Virulome (Stx type, eae/ehxA/espP, toxB, katP, tccP, espI) & VF & Yes (all patterns) \\
C7 & Stx phage type (1a/1a/2a) chromosomal not plasmid & Genomic & Yes (type); coords deferred \\
C8 & F-type virulence + separate resistance plasmid & Plasmid & Yes (PlasmidFinder) \\
C9 & Long-read recovers genes short-read misses & Methodological & Not re-run (§7 in REPORT.md) \\
\bottomrule
\end{longtable}

\section{Method}
Independent re-derivation from deposited complete genomes: NCBI efetch of 7 CP replicons;
BLAST+ per-replicon dbs vs \textbf{abricate} nucleotide reference DBs (ResFinder 3{,}206
seqs, VFDB 4{,}592, ecoli\_vf/EcOH 2{,}701, PlasmidFinder 488) at $\geq 80\%$ identity $\geq 80\%$
query coverage; MLST via PubMLST \textit{Escherichia} seqdef Achtman scheme (7 loci,
100\% full-length allele match required); custom \texttt{genome\_stats.py},
\texttt{run\_blast.py}, \texttt{summarize.py}, \texttt{mlst.py}, \texttt{stx\_location.py},
\texttt{judge.py}. Verdict via LLM-judge (free Argo \texttt{argo:gpt-5.2}) --- no regex scoring.
Wall-clock $\sim$5 min laptop CPU.

\section{Results vs paper}
\subsection*{Genome + plasmid architecture (C1--C3)}
7/7 replicons match: chromosome sizes 5{,}689{,}156 / 5{,}592{,}581 / 5{,}436{,}079 bp;
plasmids 88{,}848 / (96{,}016 + 73{,}152) / 157{,}534 bp.

\subsection*{MLST (C4)}
3/3 exact: CFSAN027343 and 346 both ST21 (16-4-12-16-9-7-7); CFSAN027350 ST29
(6-4-12-16-9-7-7), differing at \textit{adk} only --- SLV consistent with paper.

\subsection*{Acquired AMR (C5) --- the sharpest test}
ResFinder (id$\geq$80, cov$\geq$80): CFSAN027343 and 350 carry \textbf{ZERO} acquired AMR
genes anywhere. CFSAN027346 carries \textbf{6/6 EXACT}: \texttt{aph(3'')-Ib} 100\%,
\texttt{aph(6)-Id} 100\%, \textbf{\texttt{blaTEM-1B} 100\%}, \texttt{sul2} 100\%,
\texttt{tet(B)} 100\%, \texttt{dfrA8} 100\% --- all on the 73 kb plasmid CP037947. This is
an exact reproduction of the paper's central AMR finding, item-for-item and allele-for-allele.

\subsection*{Virulome (C6)}
All strain-dependent presence/absence patterns match paper Table 7: Stx type
(stx1a/stx1a/stx2a), \textit{eae}/\textit{ehxA}/\textit{espP} in all three, \textit{toxB}
only in 346+350, \textit{katP} only in 343+346, \textit{tccP} only in 346+350, \textit{espI}
only in 350. Plasmid-borne VFs (espP, toxB, katP, ehxA) confirmed on F-type plasmids.

\subsection*{Stx phage (C7)}
Types 3/3 (stx1a/stx1a/stx2a) chromosomal. Exact coordinates differ from paper Table 8
because paper reports MinION-assembly coordinates and we use PacBio-based CP GenBank
assemblies with different origin rotation/strand --- assembly-presentation, not biology.

\subsection*{Plasmid replicons (C8)}
PlasmidFinder: large virulence plasmids carry IncFIB(AP001918) + IncB/O/K/Z; AMR-bearing
73 kb plasmid is IncFII (best IncFII(pHN7A8) 98\%). Consistent with paper's F-type virulence +
distinct resistance-plasmid narrative.

\section{Verdict}
\textbf{REPLICATED.} Every core, data-grounded claim (C1--C8) reproduces on the deposited
complete genomes. AMR result is exact to the allele (\textit{blaTEM-1B}, \textit{dfrA8}).
MLST 3/3, virulome all patterns, plasmid architecture 7/7. LLM-judge (Argo gpt-5.2):
coverage 9/10, agreement 10/10, verdict REPLICATED.

\section{Genuine critique (evidence strength, shortcuts, unverified claims)}
\begin{itemize}[leftmargin=1.4em]
\item \textbf{C9 (methodological thesis: long-read beats short-read) was NOT re-run.} We
used the deposited PacBio-based CP assemblies as ground truth --- so we can confirm the
downstream \emph{biology} matches, but we did \textbf{not} independently re-assemble MinION
(SRR8335317, 3.5~GB) with CANU and MiSeq (SRR8335318/8333590-92) with SPAdes and diff
gene calls. The paper's central \emph{methodological} claim is thus \emph{architecturally
supported} (all AMR sits on the 73~kb plasmid, exactly the mobile repeat-flanked element
short-read assembly fragments) but not \emph{directly demonstrated} here. That is hours of
uicgpu compute.
\item \textbf{Assembly-accuracy claims (99.9\% consensus) not independently verified.} We
did not re-close the genomes from raw reads; we accepted the CP consensus. This means
sequence-accuracy claims of the paper are trusted, not tested. Only the biology built on
top of the sequence is tested.
\item \textbf{Stx phage coordinate mismatch (§4.6) is not scientific but presentation.}
Called out honestly but did not resolve --- would require re-rotating the CP assemblies to
paper's dnaA-frame or re-running PHASTER on our assemblies with the same origin choice.
\item \textbf{PHASTER prophage recount (paper Table 8 counts/sizes) was NOT run.} Only
Stx-carrying-phage type + chromosomal localization were checked.
\item \textbf{Single LLM-judge below policy ideal.} One judge (Argo gpt-5.2); policy prefers
3-judge robustness. Verdict defensible because the AMR result is exact and MLST/virulome
patterns are unambiguous, but scoring layer is thin.
\item \textbf{Reference-DB drift is a known-but-unquantified caveat.} We used current abricate
DBs (ResFinder 3{,}206 seqs), not the 2019 DB snapshot. If the current DB has grown, we
should expect our AMR count to be $\geq$ paper's; the fact that both are 6 for CFSAN027346
is either a happy coincidence or evidence that all extra genes fall below the 80\% cutoff
--- not disentangled.
\end{itemize}

\section{Open questions}
See \texttt{report/open\_questions.json}. Five open problems surfaced from re-reading the
paper: (Q1) full raw-read short-vs-long assembly comparison; (Q2) PHASTER prophage recount +
Stx-phage boundary consistency; (Q3) plasmid mobilization/host-range analysis for the AMR
plasmid; (Q4) genome-origin normalization for cross-study coordinate comparisons; (Q5)
Nanopore-only assembly quality on hard STEC repeats (Stx phage integration, TA systems).

\end{document}
